\chapter{इलेक्ट्रॉनिक प्रभाव और क्रियाशील मध्यवर्ती}\label{ch:b1:electronic-effects}

ट्राइफ़्लुओरोएथेनोइक अम्ल, \ce{CF3COOH}, एथेनोइक अम्ल, \ce{CH3COOH}, से लगभग दस हज़ार गुना
\omterm{def:b1:predominant-reaction:strong-acid}{प्रबल अम्ल} है, यद्यपि अम्लीय
O--H आबंध फ़्लुओरीन परमाणुओं से तीन आबंध दूर है। फ़ीनॉल
एथेनॉल से दस लाख गुना अधिक अम्लीय है, और बेंज़ीन वलय से बँधा अमोनिया अणु
--- ऐनिलीन --- मेथिल समूह से बँधे अमोनिया अणु की तुलना में दस लाख गुना कम क्षारकीय है।
कार्बनिक रसायन ऐसे विरोधाभासों से भरा है, और
उनमें से अधिकांश की व्याख्या दो प्रभाव करते हैं: $\sigma$ आबंधों के अनुदिश
विद्युत-ऋणात्मक परमाणुओं का खिंचाव, और $\pi$
निकायों पर इलेक्ट्रॉन युग्मों का फैलाव। यही दो प्रभाव तय करते हैं कि कौन-से आबंध टूटेंगे, कौन-से
मध्यवर्ती बनेंगे और अभिकर्मक कहाँ आक्रमण करेंगे: वे अगले अध्यायों की
\omterm{def:b1:elementary-steps:elementary-step}{अभिक्रिया क्रियाविधियों} का व्याकरण हैं।

\begin{recall}
\omterm{def:b1:periodicity:electronegativity}{विद्युत-ऋणात्मकता} और उसकी प्रवृत्तियाँ (\cref{ch:b1:periodicity}); \omterm{def:b1:lewis-resonance-vsepr:lewis}{लुइस
संरचनाएँ}, \omterm{def:b1:lewis-resonance-vsepr:formal-charge}{औपचारिक आवेश} और अनुनाद (\cref{ch:b1:lewis-resonance-vsepr});
$\mathrm{p}K_a$ तथा अम्लों और क्षारकों की प्रबलता
(\cref{ch:b1:predominant-reaction})। विद्यालय के खंड (कक्षा 12) ने
\omterm{def:b1:electronic-effects:nucleophile}{नाभिकरागी} से \omterm{def:b1:electronic-effects:nucleophile}{इलेक्ट्रॉनरागी} की ओर \omterm{def:b1:electronic-effects:curly-arrow}{वक्र तीर} खींचे थे।
\end{recall}

\section{प्रेरणिक प्रभाव}

\begin{definition}[प्रेरणिक प्रभाव]\label{def:b1:electronic-effects:inductive}
\emph{प्रेरणिक प्रभाव}\index{प्रेरणिक प्रभाव} $\sigma$ आबंधों के अनुदिश किसी
विद्युत-ऋणात्मक परमाणु या समूह की ओर इलेक्ट्रॉन घनत्व का विस्थापन है।
आकर्षी समूह (F, Cl, O, \ce{NO2}, \ce{CF3}) का प्रभाव $-I$ होता है;
ऐल्किल समूह और ऋणावेशित परमाणु इलेक्ट्रॉनों को धकेलते हैं, जो $+I$
प्रभाव है। बीच के आबंधों की संख्या के साथ यह प्रभाव तेज़ी से क्षीण होता है।
\end{definition}

\begin{proposition}[प्रेरणिक प्रभाव और अम्लता]\label{prop:b1:electronic-effects:inductive-pka}
कार्बोक्सिलिक अम्ल के अम्लीय स्थल के पास का आकर्षी समूह
कार्बोक्सिलेट आयन के ऋण आवेश को फैलाकर उसे स्थायी बनाता है, और
$\mathrm{p}K_a$ घटा देता है: आकर्षी समूह जितने अधिक और जितने निकट, अम्ल
उतना ही प्रबल।
\end{proposition}

\begin{proof}
$\mathrm{p}K_a$, $\mathrm{RCOOH} \rightleftharpoons
\mathrm{RCOO^-} + \ce{H+}$ की स्थिति मापता है। जो समूह कार्बोक्सिलेट से
इलेक्ट्रॉन घनत्व दूर खींचता है, वह उसके आवेश को अधिक परमाणुओं पर फैला देता है, जिससे
उसकी ऊर्जा उदासीन अम्ल की ऊर्जा से अधिक घटती है; साम्य दाईं ओर
खिसकता है और $K_a$ बढ़ता है। मापी गई शृंखला इसकी पुष्टि करती है।
\end{proof}

\begin{omfigure}
\begin{tikzpicture}
\begin{axis}[
  xbar, width=0.78\linewidth, height=5.2cm, xmin=0, xmax=5.4,
  symbolic y coords={TFA,TCA,DCA,MCA,AA}, ytick=data,
  yticklabels={\ce{CF3COOH}, \ce{CCl3COOH}, \ce{CHCl2COOH}, \ce{CH2ClCOOH}, \ce{CH3COOH}},
  xlabel={\qty{25}{\celsius} पर $\mathrm{p}K_a$}, bar width=9pt,
  nodes near coords, nodes near coords align={horizontal},
  every node near coord/.append style={font=\scriptsize},
  tick label style={font=\small}, label style={font=\small}, enlarge y limits=0.15]
\addplot[fill=omDef!45, draw=omDef] coordinates {(0.52,TFA) (0.51,TCA) (1.35,DCA) (2.87,MCA) (4.76,AA)};
\end{axis}
\end{tikzpicture}

{\small हैलोजनित एथेनोइक अम्ल। हर क्लोरीन
$\mathrm{p}K_a$ घटाता है; तीन क्लोरीनों के साथ अम्ल एथेनोइक अम्ल से लगभग दस हज़ार गुना
प्रबल हो जाता है। ट्राइफ़्लुओरो- और ट्राइक्लोरोएथेनोइक अम्ल,
ऐसे \omterm{def:b1:predominant-reaction:strong-acid}{प्रबल अम्लों} पर मापनों की अनिश्चितता की सीमा में, समान रूप से
प्रबल हैं।}
% ledger: pka:ethanoic, pka:chloroethanoic, pka:dichloroethanoic, pka:trichloroethanoic, pka:trifluoroethanoic
\end{omfigure}

ऐल्किल समूह थोड़ा-सा उलटा कार्य करते हैं: प्रोपेनोइक अम्ल ($\mathrm{p}K_a
= 4.89$) एथेनोइक अम्ल (4.76) से थोड़ा दुर्बल है, और मेथेनोइक अम्ल
(3.74), जिसमें कोई ऐल्किल समूह नहीं है, दोनों से प्रबल है।
% ledger: pka:propanoic, pka:methanoic

\section{मेसोमेरी प्रभाव}

जब \omterm{def:b1:lewis-resonance-vsepr:lewis}{एकाकी युग्म} या आवेश वहन करने वाला कोई परमाणु किसी $\pi$
निकाय से आबंधित होता है, तो उसके इलेक्ट्रॉन उस निकाय के साथ साझा हो सकते हैं; \omterm{def:b1:lewis-resonance-vsepr:resonance}{अनुनादी
संरचनाएँ} (\cref{ch:b1:lewis-resonance-vsepr}) इस साझेदारी का वर्णन करती हैं।

\begin{definition}[मेसोमेरी प्रभाव, दाता और ग्राही समूह]\label{def:b1:electronic-effects:mesomeric}
\emph{मेसोमेरी प्रभाव}\index{मेसोमेरी प्रभाव} किसी संयुग्मित $\pi$ निकाय के माध्यम से
इलेक्ट्रॉन घनत्व का विस्थापन है, जिसे \omterm{def:b1:lewis-resonance-vsepr:resonance}{अनुनादी
संरचनाएँ} दर्शाती हैं। \emph{दाता समूह}\index{दाता समूह} निकाय को
एक \omterm{def:b1:lewis-resonance-vsepr:lewis}{एकाकी युग्म} देता है, जो $+M$ प्रभाव है (\ce{-OH}, \ce{-O^-}, \ce{-NH2}, \ce{-OR}); और
\emph{ग्राही समूह}\index{ग्राही समूह} किसी विद्युत-ऋणात्मक परमाणु से बने
$\pi$ आबंध के माध्यम से निकाय से इलेक्ट्रॉन खींचता है, जो $-M$ प्रभाव है (\ce{-NO2},
\ce{-C=O}, \ce{-C#N})।
\end{definition}

\begin{omfigure}
\setchemfig{atom sep=2em}
\schemestart
\chemfig{CH_3-C(=[1]O)-[7]O^{-}}
\arrow{<->}
\chemfig{CH_3-C(-[1]O^{-})=[7]O}
\schemestop
\hspace{2.5em}
\schemestart
\chemfig{*6(=-=(-O^{-})-=-)}
\arrow{<->}
\chemfig{*6((=O)-=-\chemabove{}{\scriptstyle\ominus}-=-)}
\schemestop

\medskip
\schemestart
\chemfig{H_3C-C(=[1]O)-[7]NH_2}
\arrow{<->}
\chemfig{H_3C-C(-[1]O^{-})=[7]NH_2^{+}}
\schemestop

{\small एथेनोएट आयन (आवेश दो तुल्य ऑक्सीजनों में
साझा), फ़ीनॉक्साइड आयन (आवेश वलय में फैला हुआ; वलय वाली तीन संरचनाओं में से
एक दिखाई गई है), और एक ऐमाइड (नाइट्रोजन का \omterm{def:b1:lewis-resonance-vsepr:lewis}{एकाकी युग्म} C=O के साथ साझा, जो
ऐमाइडों को \omterm{def:b1:predominant-reaction:strong-acid}{दुर्बल क्षारक} बनाता है) की \omterm{def:b1:lewis-resonance-vsepr:resonance}{अनुनादी संरचनाएँ}।}
\end{omfigure}

\begin{proposition}[मेसोमेरी प्रभाव और अम्लता]\label{prop:b1:electronic-effects:mesomeric-pka}
जिस अम्ल का संयुग्मी क्षारक अनुनाद द्वारा अपना आवेश फैलाता है, वह उस अम्ल से प्रबल
होता है जिसका क्षारक ऐसा नहीं कर पाता: कार्बोक्सिलिक अम्ल ($\mathrm{p}K_a$ लगभग 4--5)
ऐल्कोहॉलों (लगभग 16) से कहीं अधिक प्रबल हैं, और फ़ीनॉल (10.0) बीच में आता है।
फ़ीनॉल के OH के सापेक्ष \textit{ऑर्थो} या \textit{पैरा} स्थिति में \omterm{def:b1:electronic-effects:mesomeric}{ग्राही समूह}
उसे \textit{मेटा} स्थिति की तुलना में अधिक प्रबल बनाता है।
\end{proposition}

\begin{proof}
ऐल्कॉक्साइड में आवेश एक ही ऑक्सीजन पर रहता है; कार्बोक्सिलेट में वह
दो ऑक्सीजनों में साझा होता है; फ़ीनॉक्साइड में वह ऑक्सीजन और वलय के तीन
कार्बनों में साझा होता है, जो ऑक्सीजन से कम विद्युत-ऋणात्मक हैं, इसलिए लाभ कम है।
\textit{ऑर्थो} या \textit{पैरा} कार्बन पर स्थित \ce{NO2} समूह
\ce{C=N+(-O^-)-O^-} वाली एक \omterm{def:b1:lewis-resonance-vsepr:resonance}{अनुनादी संरचना} के माध्यम से आवेश स्वयं ले सकता है;
\textit{मेटा} पर कोई भी \omterm{def:b1:lewis-resonance-vsepr:resonance}{अनुनादी संरचना} आवेश को \ce{NO2} वहन करने वाले कार्बन पर
नहीं लाती, और केवल \omterm{def:b1:electronic-effects:inductive}{प्रेरणिक प्रभाव} बचता है। मापे गए
मान: 2-नाइट्रोफ़ीनॉल 7.23, 4-नाइट्रोफ़ीनॉल 7.16, 3-नाइट्रोफ़ीनॉल 8.36,
जबकि फ़ीनॉल के लिए 9.99 और एथेनॉल के लिए 15.9।
\end{proof}
% ledger: pka:ethanol, pka:phenol, pka:2-nitrophenol, pka:3-nitrophenol, pka:4-nitrophenol

यही तर्क क्षारकों पर लागू होता है। मेथिलऐमीन (उसके
अमोनियम का $\mathrm{p}K_a$ 10.66) अमोनिया (9.25, \cref{ch:b1:predominant-reaction}) से \omterm{def:b1:predominant-reaction:strong-acid}{प्रबल क्षारक} है,
क्योंकि मेथिल समूह इलेक्ट्रॉनों को नाइट्रोजन पर धकेलता है; ऐनिलीन में \omterm{def:b1:lewis-resonance-vsepr:lewis}{एकाकी युग्म}
वलय के साथ साझा होता है और अमोनियम कहीं अधिक अम्लीय है (4.60)।
पिरिडीन (5.23) अपना \omterm{def:b1:lewis-resonance-vsepr:lewis}{एकाकी युग्म} रखता है, पर एक समतलीय वलय में स्थित ऐसे नाइट्रोजन पर
जिसमें \textit{s} लक्षण अधिक है, इसलिए वह ऐमीन की तुलना में कम उपलब्ध है।
% ledger: pka:methylammonium, pka:anilinium, pka:pyridinium, dfg:NH4+_ao

\section{आबंध का टूटना: क्रियाशील मध्यवर्ती}

\begin{definition}[समांश और विषमांश विदलन]\label{def:b1:electronic-effects:cleavage}
\omterm{def:b1:lewis-resonance-vsepr:lewis}{सहसंयोजी आबंध} \emph{समांश विदलन}\index{समांश विदलन} से टूटता है जब हर परमाणु
दो आबंधी इलेक्ट्रॉनों में से एक रखता है, जिससे दो \omterm{def:b1:electronic-effects:intermediates}{मुक्त मूलक} बनते हैं, और
\emph{विषमांश विदलन}\index{विषमांश विदलन} से जब एक परमाणु दोनों रखता है, जिससे एक
धनायन और एक ऋणायन बनते हैं।
\end{definition}

\begin{definition}[क्रियाशील मध्यवर्ती]\label{def:b1:electronic-effects:intermediates}
\emph{क्रियाशील मध्यवर्ती}\index{क्रियाशील मध्यवर्ती} ऐसी स्पीशीज़ है जो
किसी क्रियाविधि के एक चरण में बनती है और बाद के किसी चरण में खप जाती है, और इतनी
क्रियाशील होती है कि संचित नहीं होती (\cref{ch:b1:elementary-steps})।
\emph{कार्बधनायन}\index{कार्बधनायन} में एक ऐसा कार्बन होता है जिसके तीन आबंध और एक
धन आवेश हों (छह \omterm{def:b1:quantum-numbers:configuration}{संयोजकता इलेक्ट्रॉन});
\emph{कार्बऋणायन}\index{कार्बऋणायन} में ऐसा कार्बन जिसके तीन आबंध, एक \omterm{def:b1:lewis-resonance-vsepr:lewis}{एकाकी युग्म}
और एक ऋण आवेश हों; \emph{मुक्त मूलक}\index{मुक्त मूलक} में एक
\omterm{def:b1:quantum-numbers:unpaired}{अयुग्मित इलेक्ट्रॉन} वाला परमाणु होता है।
\end{definition}

\begin{definition}[कार्बन की श्रेणी]\label{def:b1:electronic-effects:carbon-class}
एक अन्य कार्बन से आबंधित कार्बन \emph{प्राथमिक
कार्बन}\index{प्राथमिक कार्बन} है, दो से आबंधित \emph{द्वितीयक
कार्बन}\index{द्वितीयक कार्बन}, तीन से आबंधित \emph{तृतीयक
कार्बन}\index{तृतीयक कार्बन}। \omterm{def:b1:electronic-effects:intermediates}{कार्बधनायन}, \omterm{def:b1:electronic-effects:intermediates}{मुक्त मूलक} या
हैलोऐल्केन संबंधित कार्बन की श्रेणी ग्रहण करता है।
\end{definition}

\begin{omfigure}
\begin{tikzpicture}[font=\small]
  % carbocation: planar, empty p orbital
  \begin{scope}
    \fill[omProp!15, draw=omProp] (0,0.15) ellipse (0.18 and 0.75);
    \fill[omProp!15, draw=omProp] (0,-0.15) ellipse (0.18 and 0.75);
    \draw[thick] (0,0) -- (-1.0,-0.25) node[left] {R};
    \draw[thick] (0,0) -- (0.95,-0.35) node[right] {R};
    \draw[thick] (0,0) -- (0.35,0.45) node[above right] {R};
    \node[fill=white, inner sep=1pt] at (0,0) {\ce{C+}};
    \node at (0,-1.75) {कार्बधनायन: समतलीय,};
    \node at (0,-2.15) {रिक्त \textit{p} कक्षक};
  \end{scope}
  % radical: nearly planar, one electron
  \begin{scope}[xshift=4.6cm]
    \draw[omMeth] (0,0.15) ellipse (0.18 and 0.75);
    \draw[omMeth] (0,-0.15) ellipse (0.18 and 0.75);
    \fill (0,0.55) circle (0.05);
    \draw[thick] (0,0) -- (-1.0,-0.25) node[left] {R};
    \draw[thick] (0,0) -- (0.95,-0.35) node[right] {R};
    \draw[thick] (0,0) -- (0.35,0.45) node[above right] {R};
    \node[fill=white, inner sep=1pt] at (0,0) {C};
    \node at (0,-1.75) {मुक्त मूलक: लगभग समतलीय,};
    \node at (0,-2.15) {एक इलेक्ट्रॉन};
  \end{scope}
  % carbanion: pyramidal with a lone pair
  \begin{scope}[xshift=9.2cm]
    \draw[thick] (0,0) -- (-0.95,-0.55) node[left] {R};
    \draw[thick] (0,0) -- (0.95,-0.55) node[right] {R};
    \draw[thick] (0,0) -- (0.2,-0.9) node[below] {R};
    \fill[omDef!15, draw=omDef] (0,0.55) ellipse (0.22 and 0.45);
    \fill (-0.07,0.65) circle (0.045); \fill (0.07,0.65) circle (0.045);
    \node[fill=white, inner sep=1pt] at (0,0) {\ce{C-}};
    \node at (0,-1.75) {कार्बऋणायन: पिरामिडी,};
    \node at (0,-2.15) {एकाकी युग्म};
  \end{scope}
\end{tikzpicture}

{\small कार्बन के तीन \omterm{def:b1:electronic-effects:intermediates}{क्रियाशील मध्यवर्तियों} की ज्यामितियाँ (R: H या
कोई ऐल्किल समूह)। चपटे \omterm{def:b1:electronic-effects:intermediates}{कार्बधनायन} पर किसी भी फलक से आक्रमण हो सकता है, जो
\cref{ch:b1:nucleophilic-substitution} के त्रिविम रसायन के लिए
महत्त्वपूर्ण बात है।}
\end{omfigure}

\begin{proposition}[कार्बधनायनों का स्थायित्व]\label{prop:b1:electronic-effects:carbocation-order}
\omterm{def:b1:electronic-effects:intermediates}{कार्बधनायन} उन समूहों द्वारा स्थायी होते हैं जो उन्हें इलेक्ट्रॉन देते हैं: तृतीयक
$>$ द्वितीयक $>$ प्राथमिक $>$ मेथिल; द्वि-आबंध (ऐलिलिक) या बेंज़ीन वलय
(बेंज़िलिक) के पास का \omterm{def:b1:electronic-effects:intermediates}{कार्बधनायन} अनुनाद द्वारा और अधिक स्थायी होता है,
और \omterm{def:b1:lewis-resonance-vsepr:lewis}{एकाकी युग्म} वाले परमाणु के पास का (\ce{R-O-CH2+}) उससे भी
अधिक।
\end{proposition}

\begin{proof}
हर ऐल्किल समूह इलेक्ट्रॉन घनत्व को इलेक्ट्रॉन-न्यून कार्बन की ओर धकेलता है
($+I$ प्रभाव, और उसके C--H आबंधों की रिक्त कक्षक के साथ एक संबंधित अन्योन्यक्रिया,
जिसका वर्णन द्वितीय वर्ष के खंड में है)। ऐलिलिक धनायन की
दो \omterm{def:b1:lewis-resonance-vsepr:resonance}{अनुनादी संरचनाएँ} होती हैं जो आवेश को किसी भी सिरे के कार्बन पर रखती हैं;
बेंज़िलिक धनायन की चार, जिनमें तीन वलय में; ऑक्सीजन का \omterm{def:b1:lewis-resonance-vsepr:lewis}{एकाकी युग्म}
धनायनी कार्बन से चौथा आबंध बनाता है, \ce{R-O+=CH2}, जिसमें हर परमाणु का
अष्टक पूरा है।
\end{proof}

\omterm{def:b1:electronic-effects:intermediates}{कार्बऋणायन} उलटे क्रम का पालन करते हैं (मेथिल $>$ प्राथमिक $>$ द्वितीयक
$>$ तृतीयक) और \omterm{def:b1:electronic-effects:mesomeric}{ग्राही समूहों} द्वारा स्थायी होते हैं; \omterm{def:b1:electronic-effects:intermediates}{मुक्त मूलक}
\omterm{def:b1:electronic-effects:intermediates}{कार्बधनायनों} वाले क्रम का ही पालन करते हैं, कम स्पष्ट रूप से।

\section{वक्र तीर, नाभिकरागी और इलेक्ट्रॉनरागी}

\begin{definition}[वक्र तीर]\label{def:b1:electronic-effects:curly-arrow}
क्रियाविधि में \emph{वक्र तीर}\index{वक्र तीर} एक इलेक्ट्रॉन युग्म की
गति दर्शाता है: वह किसी \omterm{def:b1:lewis-resonance-vsepr:lewis}{एकाकी युग्म} या आबंध से आरंभ होता है और किसी परमाणु पर
(जहाँ \omterm{def:b1:lewis-resonance-vsepr:lewis}{एकाकी युग्म} या नया आबंध बनता है) या दो परमाणुओं के बीच (नया
आबंध) समाप्त होता है। \emph{अर्ध-शीर्ष तीर}\index{अर्ध-शीर्ष तीर} \omterm{def:b1:electronic-effects:intermediates}{मुक्त मूलक} अभिक्रियाओं में
एक अकेले इलेक्ट्रॉन की गति दर्शाता है।
\end{definition}

\begin{method}[वक्र तीर खींचना]\label{met:b1:electronic-effects:curly-arrows}
\begin{enumerate}
  \item इलेक्ट्रॉनों से आरंभ कीजिए: \omterm{def:b1:lewis-resonance-vsepr:lewis}{एकाकी युग्म}, $\pi$ आबंध या $\sigma$
        आबंध, कभी धन आवेश से या इलेक्ट्रॉन-रहित परमाणु से नहीं।
  \item वहाँ समाप्त कीजिए जहाँ युग्म जाता है: दोनों परमाणुओं के बीच नया आबंध, या
        आबंध टूटने पर अधिक विद्युत-ऋणात्मक परमाणु पर \omterm{def:b1:lewis-resonance-vsepr:lewis}{एकाकी युग्म}।
  \item हर चरण के बाद आवेश गिनिए: जो परमाणु नए आबंध को \omterm{def:b1:lewis-resonance-vsepr:lewis}{एकाकी युग्म} देता है, वह
        $+1$ पाता है, और जो परमाणु युग्म को \omterm{def:b1:lewis-resonance-vsepr:lewis}{एकाकी युग्म} के रूप में ग्रहण करता है,
        वह $-1$ पाता है; कुल आवेश संरक्षित रहता है।
  \item C, N, O पर कभी अष्टक से अधिक न होने दें: जब एक युग्म आता है, तो दूसरे को
        जाना ही होगा।
\end{enumerate}
\end{method}

\begin{omfigure}
\setchemfig{atom sep=2.2em}
\schemestart
\chemfig{H_3C-C(-[2]H)(-[6]H)-[@{b}]@{br}Br}
\arrow{->[विषमांश विदलन]}[,1.5]
\chemfig{H_3C-C(-[2]H)(-[6]H)^{+}}\+\chemfig{Br^{-}}
\schemestop
\chemmove{\draw[->, shorten <=2pt, shorten >=2pt](b) .. controls +(90:6mm) and +(90:6mm) .. (br);}

\medskip
\schemestart
\chemfig{H@{o}O^{-}}\+\chemfig{@{h}H-@{oc}O-H}
\arrow{<=>}
\chemfig{H_2O}\+\chemfig{HO^{-}}
\schemestop
\chemmove{\draw[->, thick, shorten <=2pt, shorten >=2pt](o) .. controls +(-80:8mm) and +(-100:8mm) .. (h);}
\par\vspace{5mm}

{\small \omterm{def:b1:electronic-effects:curly-arrow}{वक्र तीर}। ऊपर: C--Br आबंध का \omterm{def:b1:electronic-effects:cleavage}{विषमांश विदलन}, युग्म
ब्रोमीन के पास जाता है। नीचे: जल से हाइड्रॉक्साइड को प्रोटॉन का स्थानांतरण, जिसे
हाइड्रोजन पर क्षारक के \omterm{def:b1:lewis-resonance-vsepr:lewis}{एकाकी युग्म} के आक्रमण के रूप में लिखा गया है (तब
जल का O--H आबंध टूटता है, उसका युग्म ऑक्सीजन पर रहता है; दूसरा तीर
\cref{exo:b1:electronic-effects:4} के लिए छोड़ दिया गया है)।}
\end{omfigure}

\begin{definition}[नाभिकरागी, इलेक्ट्रॉनरागी, निर्गामी समूह]\label{def:b1:electronic-effects:nucleophile}
\emph{नाभिकरागी}\index{नाभिकरागी} वह स्पीशीज़ है जो आबंध बनाने के लिए
हाइड्रोजन के अतिरिक्त किसी परमाणु को इलेक्ट्रॉन युग्म देती है (एक \omterm{def:b1:lewis-resonance-vsepr:lewis-acid}{लुइस
क्षारक}); \emph{इलेक्ट्रॉनरागी}\index{इलेक्ट्रॉनरागी} वह स्पीशीज़ है जो
ऐसा युग्म ग्रहण करती है (एक \omterm{def:b1:lewis-resonance-vsepr:lewis-acid}{लुइस अम्ल}, प्रायः आंशिक
धन आवेश वहन करने वाला कार्बन)। किसी स्पीशीज़ की \emph{नाभिकरागिता}\index{नाभिकरागिता}
नाभिकरागी के रूप में उसकी अभिक्रियाशीलता है, जिसे \omterm{def:b1:rate-laws:order}{दर स्थिरांकों} से मापा जाता है।
\emph{निर्गामी समूह}\index{निर्गामी समूह} वह समूह है जो
कार्बन से बना आबंध विषमांशतः टूटने पर \omterm{def:b1:lewis-resonance-vsepr:lewis}{आबंधी युग्म} लेकर निकल जाता है।
\end{definition}

\begin{proposition}[नाभिकरागिता और क्षारकता]\label{prop:b1:electronic-effects:nucleophilicity-basicity}
एक ही परमाणु से आक्रमण करने वाले \omterm{def:b1:electronic-effects:nucleophile}{नाभिकरागियों} के लिए \omterm{def:b1:electronic-effects:nucleophile}{नाभिकरागिता}
क्षारकता का अनुसरण करती है: \ce{CH3O-} $>$ \ce{OH-} $>$ \ce{CH3COO-} $>$ \ce{H2O}।
\omterm{def:b1:intermolecular-forces-solvents:solvent-classes}{प्रोटिक विलायकों} में आवर्त सारणी के किसी समूह में नीचे जाने पर यह नियम विफल हो जाता है, जहाँ
बड़े, ध्रुवणीय, दुर्बल रूप से विलायकित \omterm{def:b1:electronic-effects:nucleophile}{नाभिकरागी} \omterm{def:b1:predominant-reaction:strong-acid}{दुर्बल क्षारक} होते हुए भी
तेज़ी से अभिक्रिया करते हैं (\ce{I-} $>$ \ce{Br-} $>$ \ce{Cl-} $>$ \ce{F-}), और
भारी-भरकम क्षारकों के लिए भी, जो दुर्बल \omterm{def:b1:electronic-effects:nucleophile}{नाभिकरागी} होते हैं।
\end{proposition}

\begin{proof}
क्षारकता प्रोटॉन पर आक्रमण के साम्य को मापती है,
\omterm{def:b1:electronic-effects:nucleophile}{नाभिकरागिता} कार्बन पर आक्रमण की दर को; दोनों इलेक्ट्रॉन युग्म की
उपलब्धता के साथ बढ़ती हैं, इसलिए एक ही परमाणु के भीतर समानांतर चलती हैं।
दर इस पर भी निर्भर करती है कि \omterm{def:b1:electronic-effects:nucleophile}{नाभिकरागी} से विलायक हटाने के लिए कितनी ऊर्जा चाहिए
और भीड़ वाले कार्बन की ओर उसके इलेक्ट्रॉन मेघ का
विरूपण कितना है: \ce{F-} जैसे छोटे ऋणायन \omterm{def:b1:intermolecular-forces-solvents:solvent-classes}{प्रोटिक विलायक} अणुओं से
प्रबलता से बँधे रहते हैं (\cref{ch:b1:intermolecular-forces-solvents}), और
भारी-भरकम समूह कार्बन तक पहुँच में बाधा डालते हैं, पर प्रोटॉन तक नहीं।
\end{proof}

अच्छा \omterm{def:b1:electronic-effects:nucleophile}{निर्गामी समूह} एक \omterm{def:b1:predominant-reaction:strong-acid}{दुर्बल क्षारक} होता है, जो युग्म ले जाने के बाद स्थायी रहता है:
\ce{I-}, \ce{Br-}, \ce{Cl-}, जल, सल्फ़ोनेट; \ce{OH-} और \ce{RO-}, जो \omterm{def:b1:predominant-reaction:strong-acid}{प्रबल
क्षारक} हैं, खराब \omterm{def:b1:electronic-effects:nucleophile}{निर्गामी समूह} हैं --- यही कारण है कि ऐल्कोहॉलों को
प्रतिस्थापन से पहले सक्रियित करना पड़ता है (\cref{ch:b1:alcohol-activation})।

\section{अभ्यास}

\begin{exercise}[$\star$]\label{exo:b1:electronic-effects:1}
2-मेथिलब्यूटेन और 2-ब्रोमो-2-मेथिलप्रोपेन के हर कार्बन को
प्राथमिक, द्वितीयक, तृतीयक या चतुष्क (चार कार्बनों से आबंधित) में वर्गीकृत कीजिए।
सूत्र \ce{C4H9Br} वाले हर ब्रोमाइड से कौन-सा \omterm{def:b1:electronic-effects:intermediates}{कार्बधनायन} सबसे आसानी से बनता है?
\end{exercise}

\begin{exercise}[$\star$]\label{exo:b1:electronic-effects:2}
$\mathrm{p}K_a$ मानों से क्लोरोएथेनोइक और एथेनोइक अम्लों के, तथा ट्राइफ़्लुओरोएथेनोइक
और एथेनोइक अम्लों के \omterm{def:b1:predominant-reaction:acidity-constant}{अम्लता स्थिरांकों} का अनुपात
परिकलित कीजिए।
\end{exercise}

\begin{exercise}[$\star$]\label{exo:b1:electronic-effects:3}
एथेनोएट आयन, 4-नाइट्रोफ़ीनॉक्साइड
आयन (आवेश नाइट्रो समूह के एक ऑक्सीजन पर रखते हुए) और ऐलिल
धनायन \ce{CH2=CH-CH2+} की \omterm{def:b1:lewis-resonance-vsepr:resonance}{अनुनादी संरचनाएँ} खींचिए।
\end{exercise}

\begin{exercise}[$\star$]\label{exo:b1:electronic-effects:4}
चित्र के प्रोटॉन स्थानांतरण के और
\ce{NH3 + H3O+ -> NH4+ + H2O} के \omterm{def:b1:electronic-effects:curly-arrow}{वक्र तीर} पूरे कीजिए। आवेशों की जाँच कीजिए।
\end{exercise}

\begin{exercise}[$\star\star$]\label{exo:b1:electronic-effects:5}
बढ़ती अम्लता के क्रम में रखिए: एथेनॉल, फ़ीनॉल, एथेनोइक अम्ल,
4-नाइट्रोफ़ीनॉल, ट्राइक्लोरोएथेनोइक अम्ल, जल (युग्म \ce{H2O}/\ce{OH-} का $\mathrm{p}K_a$
14.0)। हर चरण का औचित्य किसी प्रेरणिक या
मेसोमेरी तर्क से दीजिए।
\end{exercise}

\begin{exercise}[$\star\star$]\label{exo:b1:electronic-effects:6}
बढ़ती क्षारकता के क्रम में रखिए: ऐनिलीन, मेथिलऐमीन, अमोनिया, पिरिडीन।
एक \omterm{def:b1:lewis-resonance-vsepr:resonance}{अनुनादी संरचना} से ऐनिलीन की स्थिति समझाइए।
\end{exercise}

\begin{exercise}[$\star\star$]\label{exo:b1:electronic-effects:7}
3-नाइट्रोफ़ीनॉल 4-नाइट्रोफ़ीनॉल से दुर्बल पर फिर भी फ़ीनॉल से प्रबल
क्यों है? वे \omterm{def:b1:lewis-resonance-vsepr:resonance}{अनुनादी संरचनाएँ} खींचिए जो दोनों कथनों का औचित्य बताती हैं।
\end{exercise}

\begin{exercise}[$\star\star$]\label{exo:b1:electronic-effects:8}
\omterm{def:b1:electronic-effects:intermediates}{कार्बधनायनों} \ce{CH3+}, \ce{(CH3)3C+}, \ce{CH2=CH-CH2+},
\ce{CH3CH2+}, \ce{C6H5CH2+}, \ce{CH3OCH2+} को क्रम में रखिए, और औचित्य दीजिए।
\end{exercise}

\begin{exercise}[$\star\star$]\label{exo:b1:electronic-effects:9}
इनमें \omterm{def:b1:electronic-effects:nucleophile}{नाभिकरागी} और \omterm{def:b1:electronic-effects:nucleophile}{इलेक्ट्रॉनरागी} पहचानिए: \ce{CH3Br + OH-};
\ce{H2O + H+}; सायनाइड आयन के साथ एथेनल; \ce{BF3 + NH3}। \omterm{def:b1:electronic-effects:curly-arrow}{वक्र
तीर} खींचिए।
\end{exercise}

\begin{exercise}[$\star\star\star$]\label{exo:b1:electronic-effects:10}
एक विलयन में \qty{0.010}{mol/L} ट्राइक्लोरोएथेनोइक अम्ल है। क्या
\omterm{def:b1:predominant-reaction:strong-acid}{दुर्बल अम्ल} का सूत्र $\mathrm{pH} = \frac12(\mathrm{p}K_a - \log C)$ मान्य है?
pH ठीक से परिकलित कीजिए, और \omterm{def:b1:predominant-reaction:dissociation}{वियोजन मात्रा} भी।
\end{exercise}

\begin{exercise}[$\star\star\star$]\label{exo:b1:electronic-effects:11}
जल में मापे गए ऐल्कोहॉलों के $\mathrm{p}K_a$ मेथेनॉल
(15.5) से एथेनॉल (15.9), प्रोपेन-2-ऑल (17.1) और 2-मेथिलप्रोपेन-2-ऑल
(19.2) तक बढ़ते हैं। दो व्याख्याएँ दीजिए, एक ऐल्कॉक्साइड पर ऐल्किल समूहों के
\omterm{def:b1:electronic-effects:inductive}{प्रेरणिक प्रभाव} के माध्यम से, और एक ऐल्कॉक्साइड के \omterm{def:b1:intermolecular-forces-solvents:dissolution}{विलायकन} के माध्यम से।
% ledger: pka:methanol, pka:ethanol, pka:2-propanol, pka:tBuOH
\end{exercise}

\begin{exercise}[$\star\star\star$]\label{exo:b1:electronic-effects:12}
ऐमाइड की \omterm{def:b1:lewis-resonance-vsepr:resonance}{अनुनादी संरचना} का उपयोग करके समझाइए कि ऐमाइड न तो नाइट्रोजन पर क्षारकीय हैं
न \omterm{def:b1:electronic-effects:nucleophile}{नाभिकरागी}, जबकि ऐमीन दोनों हैं।
\omterm{def:b1:predominant-reaction:strong-acid}{प्रबल अम्ल} में ऐमाइड का कौन-सा परमाणु प्रोटॉनित होता है?
\end{exercise}

\section{समस्या: ट्राइफ़्लुओरोएथेनोइक अम्ल इतना प्रबल क्यों है}

\begin{problem}[{सप्ताहांत समस्या --- हैलोजनित अम्लों की
प्रेरणिक शृंखला, कार्बोक्सिलेट बनाम ऐल्कॉक्साइड, तीन
नाइट्रोफ़ीनॉल, एक प्रधान-अभिक्रिया परिकलन, और ट्राइफ़्लुओरोएथेनोइक तथा एथेनोइक
अम्लों के अम्लता स्थिरांकों का अनुपात}]
\label{pb:b1:electronic-effects:1}
आँकड़े (\qty{25}{\celsius} पर $\mathrm{p}K_a$): एथेनोइक अम्ल 4.76,
क्लोरोएथेनोइक 2.87, डाइक्लोरोएथेनोइक 1.35, ट्राइक्लोरोएथेनोइक 0.51,
ट्राइफ़्लुओरोएथेनोइक 0.52, एथेनॉल 15.9, फ़ीनॉल 9.99, 2-नाइट्रोफ़ीनॉल 7.23,
3-नाइट्रोफ़ीनॉल 8.36, 4-नाइट्रोफ़ीनॉल 7.16; $\mathrm{p}K_e = 14.00$।

\medskip
\noindent\textbf{भाग I --- प्रेरणिक शृंखला।}
\begin{enumerate}
  \item एथेनोइक से ट्राइक्लोरोएथेनोइक अम्ल तक हर चरण के लिए $K_a$ का अनुपात
        परिकलित कीजिए।
  \item क्या हर क्लोरीन का प्रभाव समान है? टिप्पणी कीजिए।
  \item \omterm{def:b1:electronic-effects:inductive}{प्रेरणिक प्रभाव} से इस प्रवृत्ति की व्याख्या कीजिए।
  \item किसी लंबी शृंखला के उस कार्बन पर क्लोरीन, जो COOH से सबसे दूर है,
        लगभग कोई प्रभाव क्यों नहीं डालेगा?
  \item फ़्लुओरीन क्लोरीन से अधिक विद्युत-ऋणात्मक है। ट्राइफ़्लुओरो- और ट्राइक्लोरो- अम्लों की
        तुलना क्या दर्शाती है, और ऐसे $\mathrm{p}K_a$ को यथार्थता से मापना
        कठिन क्यों है?
  \item pH 3 पर इनमें से हर अम्ल किस रूप में है?
\end{enumerate}

\medskip
\noindent\textbf{भाग II --- कार्बोक्सिलेट और ऐल्कॉक्साइड।}
\begin{enumerate}[resume]
  \item एथेनोएट आयन की दोनों \omterm{def:b1:lewis-resonance-vsepr:resonance}{अनुनादी संरचनाएँ} खींचिए।
  \item वे आयन की दोनों C--O आबंध लंबाइयों के बारे में क्या पूर्वानुमान देती हैं?
  \item एथॉक्साइड आयन की कोई तुल्य संरचनाएँ क्यों नहीं हैं?
  \item एथेनोइक अम्ल और एथेनॉल के $K_a$ का अनुपात परिकलित कीजिए।
  \item फ़ीनॉल को दोनों के बीच रखिए और फ़ीनॉक्साइड की \omterm{def:b1:lewis-resonance-vsepr:resonance}{अनुनादी
        संरचनाओं} से औचित्य दीजिए।
  \item एथेनॉल, फ़ीनॉल और एथेनोइक अम्ल में से कौन-से हाइड्रोजनकार्बोनेट विलयन के साथ
        लगभग पूर्णतः अभिक्रिया करते हैं (\ce{CO2,H2O}/\ce{HCO3-} के लिए
        $\mathrm{p}K_a$ 6.37)?
\end{enumerate}

\medskip
\noindent\textbf{भाग III --- नाइट्रोफ़ीनॉल।}
\begin{enumerate}[resume]
  \item तीनों नाइट्रोफ़ीनॉलों और फ़ीनॉल को अम्लता के क्रम में रखिए।
  \item 4-नाइट्रोफ़ीनॉक्साइड की एक ऐसी \omterm{def:b1:lewis-resonance-vsepr:resonance}{अनुनादी संरचना} खींचिए जिसमें आवेश
        नाइट्रो समूह पर हो।
  \item दिखाइए कि 3-नाइट्रोफ़ीनॉक्साइड के लिए ऐसी कोई संरचना नहीं है।
  \item 3-नाइट्रोफ़ीनॉल फिर भी फ़ीनॉल से अधिक अम्लीय क्यों है?
  \item 2-नाइट्रोफ़ीनॉल 4-नाइट्रोफ़ीनॉल से थोड़ा दुर्बल है, यद्यपि दोनों में
        \omterm{def:b1:electronic-effects:mesomeric}{मेसोमेरी प्रभाव} है; उसके OH और पास के नाइट्रो समूह से जुड़ी एक
        व्याख्या प्रस्तावित कीजिए।
  \item pH 7.5 पर कौन-से नाइट्रोफ़ीनॉल अधिकांशतः ऋणायन रूप में हैं? इससे
        4-नाइट्रोफ़ीनॉल \omterm{def:b1:titration-methods:indicator}{रंग सूचक} क्यों बनता है?
\end{enumerate}

\medskip
\noindent\textbf{भाग IV --- जल में ट्राइफ़्लुओरोएथेनोइक अम्ल।}
\begin{enumerate}[resume]
  \item \qty{0.010}{mol/L} पर ट्राइफ़्लुओरोएथेनोइक अम्ल का pH परिकलित कीजिए
        (दुर्बल वियोजन मत मानिए)।
  \item \qty{0.010}{mol/L} पर एथेनोइक अम्ल का pH परिकलित कीजिए।
  \item सोडियम एथेनोएट के साथ ट्राइफ़्लुओरोएथेनोइक अम्ल की अभिक्रिया लिखिए
        और उसका स्थिरांक परिकलित कीजिए।
  \item अनुपात $K_a(\ce{CF3COOH})/K_a(\ce{CH3COOH})$ बताइए, दस की घात के
        रूप में और एक संख्या के रूप में।
\end{enumerate}
\end{problem}
